GPT-3 and our analysis of it have a number of limitations.  Below we describe some of these and suggest directions for future work.

First, despite the strong quantitative and qualitative improvements of GPT-3, particularly compared to its direct predecessor GPT-2, it still has notable weaknesses in text synthesis and several NLP tasks.  On text synthesis, although the overall quality is high, GPT-3 samples still sometimes repeat themselves semantically at the document level, start to lose coherence over sufficiently long passages, contradict themselves, and occasionally contain non-sequitur sentences or paragraphs.  We will release a collection of 500 uncurated unconditional samples to help provide a better sense of GPT-3’s limitations and strengths at text synthesis. Within the domain of discrete language tasks, we have noticed informally that GPT-3 seems to have special difficulty with ``common sense physics'', despite doing well on some datasets (such as PIQA \cite{bisk2019piqa}) that test this domain.  Specifically GPT-3 has difficulty with questions of the type ``If I put cheese into the fridge, will it melt?''.  Quantitatively, GPT-3’s in-context learning performance has some notable gaps on our suite of benchmarks, as described in Section \ref{section:Results}, and in particular it does little better than chance when evaluated one-shot or even few-shot on some ``comparison'' tasks, such as determining if two words are used the same way in a sentence, or if one sentence implies another (WIC and ANLI respectively), as well as on a subset of reading comprehension tasks.  This is especially striking given GPT-3’s strong few-shot performance on many other tasks.

GPT-3 has several structural and algorithmic limitations, which could account for some of the issues above. We focused on exploring in-context learning behavior in autoregressive language models because it is straightforward to both sample and compute likelihoods with this model class. As a result our experiments do not include any bidirectional architectures or other training objectives such as denoising. This is a noticeable difference from much of the recent literature, which has documented improved fine-tuning performance when using these approaches over standard language models~\cite{raffel2019t5}. Thus our design decision comes at the cost of potentially worse performance on tasks which empirically benefit from bidirectionality. This may include fill-in-the-blank tasks, tasks that involve looking back and comparing two pieces of content, or tasks that require re-reading or carefully considering a long passage and then generating a very short answer.  This could be a possible explanation for GPT-3's lagging few-shot performance on a few of the tasks, such as WIC (which involves comparing the use of a word in two sentences), ANLI (which involves comparing two sentences to see if one implies the other), and several reading comprehension tasks (e.g. QuAC and RACE).  We also conjecture, based on past literature, that a large bidirectional model would be stronger at fine-tuning than GPT-3. Making a bidirectional model at the scale of GPT-3, and/or trying to make bidirectional models work with few- or zero-shot learning, is a promising direction for future research, and could help achieve the ``best of both worlds''.

A more fundamental limitation of the general approach described in this paper -- scaling up any LM-like model, whether autoregressive or bidirectional -- is that it may eventually run into (or could already be running into) the limits of the pretraining objective. Our current objective weights every token equally and lacks a notion of what is most important to predict and what is less important. \cite{roberts2020much} demonstrate benefits of customizing prediction to entities of interest. Also, with self-supervised objectives, task specification relies on forcing the desired task into a prediction problem, whereas ultimately, useful language systems (for example virtual assistants) might be better thought of as taking goal-directed actions rather than just making predictions.  Finally, large pretrained language models are not grounded in other domains of experience, such as video or real-world physical interaction, and thus lack a large amount of context about the world \cite{bisk2020experience}.  For all these reasons, scaling pure self-supervised prediction is likely to hit limits, and augmentation with a different approach is likely to be necessary.  Promising future directions in this vein might include learning the objective function from humans \cite{ziegler2019finetuning}, fine-tuning with reinforcement learning, or adding additional modalities such as images to provide grounding and a better model of the world \cite{chen2019uniter}.

Another limitation broadly shared by language models is poor sample efficiency during pre-training.  While GPT-3 takes a step towards test-time sample efficiency closer to that of humans (one-shot or zero-shot), it still sees much more text during pre-training than a human sees in the their lifetime \cite{linzen2020can}.  Improving pre-training sample efficiency is an important direction for future work, and might come from  grounding in the physical world to provide additional information, or from algorithmic improvements. 

A limitation, or at least uncertainty, associated with few-shot learning in GPT-3 is ambiguity about whether few-shot learning actually learns new tasks ``from scratch'' at inference time, or if it simply recognizes and identifies tasks that it has learned during training.  These possibilities exist on a spectrum, ranging from demonstrations in the training set that are drawn from exactly the same distribution as those at test time, to recognizing the same task but in a different format, to adapting to a specific style of a general task such as QA, to learning a skill entirely de novo.  Where GPT-3 is on this spectrum may also vary from task to task.  Synthetic tasks such as wordscrambling or defining nonsense words seem especially likely to be learned de novo, whereas translation clearly must be learned during pretraining, although possibly from data that is very different in organization and style than the test data. Ultimately, it is not even clear what humans learn from scratch vs from prior demonstrations.  Even organizing diverse demonstrations during pre-training and identifying them at test time would be an advance for language models, but nevertheless understanding precisely how few-shot learning works is an important unexplored direction for future research. 

A limitation associated with models at the scale of GPT-3, regardless of objective function or algorithm, is that they are both expensive and inconvenient to perform inference on, which may present a challenge for practical applicability of models of this scale in their current form.  One possible future direction to address this is distillation \cite{hinton2015distilling} of large models down to a manageable size for specific tasks. Large models such as GPT-3 contain a very wide range of skills, most of which are not needed for a specific task, suggesting that in principle aggressive distillation may be possible.  Distillation is well-explored in general \cite{liu2019improving} but has not been tried at the scale of hundred of billions parameters; new challenges and opportunities may be associated with applying it to models of this size.

Finally, GPT-3 shares some limitations common to most deep learning systems -- its decisions are not easily interpretable, it is not necessarily well-calibrated in its predictions on novel inputs as observed by the much higher variance in performance than humans on standard benchmarks, and it retains the biases of the data it has been trained on.  This last issue -- biases in the data that may lead the model to generate stereotyped or prejudiced content -- is of special concern from a societal perspective, and will be discussed along with other issues in the next section on Broader Impacts (Section \ref{section:Broader_Impacts}).

